%%%%%%%%%%%%%%
% Modelo XXXXXXXXXXXXXXX
% Autor: Leonardo Garcez Dalenogari Alba
% Ano: 2022
% Contato: leonardo.2017007421@aluno.iffar.edu.br

% IMPORTANTE: 
%   O modelo utiliza o compilador XeLaTex (talvez possam ser utilizados outros, mas não é garantido o correto funcionamento, visto que seria necessário testar, porém, eu sei que alguns vão apresentar problemas na hora de compilar, como o pdflatex)
%   Se você estiver utilizando o Overleaf, com o documento aberto, clique em Menu, no canto superior esquerdo, e selecione o XeLaTex como compilador
%   Também é utilizado o biber para gerenciar as referências bibliográficas
%%%%%%%%%%%%%%

\documentclass[
    12pt, % Tamanho de fonte definido para 12, se você precisar utilizar um texto em tamanho 10, utilize o comando \footnotesize, que será impresso em tamanho 10
    oneside,
    %twoside, % A opção permite gerar o documento utilizando o verso das páginas, porém, ela não indica à impressora que você está utilizando isso ou que ela deva imprimir nos dois lados das folhas (fonte: https://en.wikibooks.org/wiki/LaTeX/Document_Structure)
    %openany,
    %openright,
    a4paper,
    english,
    brazil
]{report}

% Pacotes que envolvem codificação do conteúdo. Essenciais.
\usepackage[T1]{fontenc}
% \usepackage[utf8]{inputenc} % Não é necessário se você estiver utilizando um compilador LaTeX baseado já em UTF-8

% Responsável pela quebra silábica de palavras no final de linha de forma correta para o português
\usepackage[brazil]{babel}
% O pacote babel recomenda o uso do csquotes, para garantir que textos citados estejam escritos corretamente de acordo com as regras do seu idioma.
\usepackage[portuguese=brazilian]{csquotes}

%% Pacotes para formatação de fonte do documento
\usepackage{fontspec}
% Pode ser Times New Roman:
% \setmainfont{Times New Roman} 
% Ou pode ser Arial, também:
\setmainfont{Arial}

% Pacote responsável por adicionar símbolos matemáticos. Símbolos, por exemplo, também necessários para o correto referenciamento de leis e decretos
\usepackage{amssymb}
% Só que também é importante realizar a configuração para adicionar os comandos que imprimem esses símbolos associados aos caracteres que imprimem, assim, você pode adicionar o símbolo diretamente no texto
\usepackage{newunicodechar} % Pacote que permite definir novos caracteres
\usepackage{config/unicode-char} % Arquivo de configuração com a listagem desses caracteres

%% Pacotes para formatação de conteúdo na página
\usepackage[left=3cm,top=3cm,right=2cm,bottom=2cm]{geometry} % Pacote responsável por formatar as dimensões das margens das páginas.
\usepackage{setspace} % Adiciona opções de espaçamento de texto
\usepackage{ragged2e} % Adiciona opções mais completas para alinhamento de texto
\usepackage{indentfirst} % Faz com que o primeiro parágrafo de cada seção possua indentação também

% Pacote que permite criar estilos para o cabeçalho
\usepackage{fancyhdr}
\usepackage{config/cabecalho} % Cabeçalho customizado para o modelo

%% Adiciono os pacotes para configurar e estilizar o sumário e as seções
% \usepackage{tocloft}
\usepackage{titletoc} % Não consegui criar pelo titletoc a parte de todas as seções estarem alinhadas, então vou usar o tocloft
\usepackage{titlesec}
\usepackage{alphalph} % Permite lidar com contadores que possuam valor maior que 26 utilizando o alfabeto (o número de letras), ou seja, se passar de 26, vai utilizar AA, AB, AC... Nativamente o LaTeX não realiza esse controle
\usepackage{config/apendice-anexo} % Criação dos comandos para apêndice e anexo

%% Agora customizo as seções e sumário para seguirem as normas do IFFar (itens adicionais do sumário, como apêndice e anexos, também são configurados nesses arquivos)
\usepackage{config/sumario} % Importa o sumário feito para o modelo
%% Adiciona a possibilidade de estilizar os títulos de seções, subseções e afins
\usepackage{config/titulo-sections} % Importa as estilizações feitas das seções e subseções para o modelo

% Pacote para facilitar a customização de listas numeradas, que serão utilizadas para as alíneas e subalíneas
\usepackage{enumitem}
\usepackage{config/enumerate}

\usepackage{float} % Adiciona configurações adicionais para elementos float (figuras, tabelas, quadros, etc.)
\usepackage{tabularx} % Tabelas no LaTeX são limitadas em algumas características, este pacote resolve algumas dessas falhas
\usepackage{graphicx} % Adiciona a capacidde de adicionar figuras, além de outras funcionalidade que podem ser utilizadas em quadros e tabelas.
% Figuras podem ser imagens (SVG, por padrão, não funciona), mas documentos PDF também são compatíveis
\graphicspath{ {./img/} } % Adiciona um caminho padrão para arquivo
% Permite modificar a apresentação das legendas (caption) de figuras, quadros e tabelas
\usepackage[labelsep=endash, font=singlespacing, justification=centering]{caption} % Modifica a separação da numeração da legenda com a descrição dela
\usepackage{newfloat} % Permite criar floats, como o caso dos quadros, de forma prática
\DeclareFloatingEnvironment[listname={Lista de quadros}, name={Quadro}]{quadro} % Crio o elemento float Quadro
% Removo a contabilização de capítulo na numeração das figuras, tabelas e quadros (isso ocorre ao utilizar o documentclass de book ou report)
\counterwithout{figure}{chapter} 
\counterwithout{table}{chapter}
\counterwithout{quadro}{chapter}
% Obs.: Se você precisar criar outro elemento (e outra lista), como de gráficos ou de qualquer outro tipo, é, teoricamente, bastante simples: 
%   - Você declara o novo ambiente com o \DeclareFloatingEnvironment, junto com o nome do elemento e o título de sua lista
%   - Remove a contabilização do capítulo com o \counterwithout (apesar que acho que isso pode ser realizado no próprio \DeclareFloatingEnvironment)
%   - E depois vai no final do arquivo config/sumario.sty e copia a estilização usada para os elementos de figura ou tabela, realizando a modificação do nome do elemento

%% Pacotes para adicionar o uso de citações e referências no padrão ABNT
% Permite customizar as notas de rodapé. Necessário para a indentação das notas
\usepackage[hang]{footmisc}
\usepackage{config/rodape} % Configurações das notas de rodapé
% Usa o biber por ser uma versão mais parruda de biblatex, que é uma versão mais parruda do bibtex, ou algo assim. É importante por resolver uma quantia de erros por conta de caracteres, principalmente na URL das referências
\usepackage[backend=biber, style=abnt]{biblatex}
\addbibresource{bibliografia.bib} % Você pode adicionar mais de um arquivo.bib, apenas adicione outra linha com o nome do segundo arquivo
\usepackage{config/bibliografia-abnt} % Ajustes necessários por causa que o pacote biblatex-abnt não foi atualizado para algumas das mais novas normas da ABNT

%Pacote responsável por permitir colocar url. Ele cria link de atalho no sumário, citações e referência de figuras, tabelas, etc.
\usepackage[hidelinks]{hyperref}
\hypersetup{
    unicode=true
}
\usepackage{url}

\usepackage{lipsum} % Pacote meramente cosmético, para testes. Remova ele quando for utilizar o modelo.

\begin{document}
\onehalfspacing
\justify

\setlength\parindent{1.25cm} % Configura o espaçamento da indentação dos parágrafos
\setlength{\parskip}{1.241pt} % Configura o espaço separação entre parágrafos.

\begin{SecoesNaoNumeradas}
  \vspace*{\fill}

% Tamanho total da página: 21cm
% Margem esquerda: 3cm
% Margem direita: 2cm
% Recuo necessário: 4cm
% Espaço restante: 12cm
\hspace{\fill}
\begin{minipage}{12cm}
    \setlength{\parindent}{1.25cm}
    
    Este trabalho é dedicado [especificar o texto da dedicatória, escrito pelo(a) autor(a), registrando a quem dedica o Trabalho Acadêmico-Científico realizado ou prestando uma homenagem a alguém que considere importante no transcorrer do seu processo formativo (pessoal e/ou acadêmico-profissional), ainda que essa pessoa não tenha contribuído diretamente para os resultados obtidos no trabalho].
\end{minipage}

% Recomendação do Guia de Normalização de iniciar cada elemento pré-textual e textual no anverso da página
\OnesideTwoside{\clearpage}{\cleardoublepage}

  \chapter*{Agradecimentos}
\thispagestyle{empty}

\lipsum[2]
  \vspace*{\fill}

% Largura da folha: 21cm
% Margem à esquerda: 3cm
% Margem à direita: 2cm
% Espaço para texto: 16cm;
% Margem à esquerda necessária: 10cm
% Se for margem de 10cm da página inteira, tamanho de: 11cm
% Se for margem de 10cm da parte de texto, sem considerar os 3cm: 6cm 
\hspace{\fill}
\begin{minipage}{6cm}
    \begin{singlespace} % A epígrafe tem que usar espaçamento símples
        % Seu texto recomenda-se em itálico e autor em negrito
        \textit{Lorem ipsum dolor sit amet, consectetur adipiscing elit. Curabitur et neque metus. Ut volutpat sem quis lacus rutrum, a viverra odio tincidunt. Quisque eu mollis enim. Suspendisse porta vehicula pharetra. Sed magna risus, consequat vitae tempus at, pulvinar ut sapien. Donec sit amet metus feugiat massa molestie venenatis id vel est. Fusce dapibus nisi vitae urna suscipit, non hendrerit mauris consequat. Nunc in quam vel dolor auctor interdum eget nec metus. Donec ornare libero euismod mi vehicula tincidunt. Pellentesque fermentum vulputate lacus, sit amet facilisis dolor faucibus et} (\textbf{Autor}).
    \end{singlespace}
\end{minipage}

% Recomendação do Guia de Normalização de iniciar cada elemento pré-textual e textual no anverso da página
\OnesideTwoside{\clearpage}{\cleardoublepage}
  \chapter*{Resumo}\myIdxResumoVernac

% O \noindent é necessário por causa que, como eu configurei para indentar o primeiro parágrafo no texto, como exigido, o LaTeX iria normalmente indentar o parágrafo
% Só escrever o teu resuminho, sem estresse
\noindent [Inserir o resumo que representa uma síntese do trabalho e cujo texto é escrito na língua vernácula, de modo objetivo, sucinto e claro. Deve ser digitado como um bloco único de parágrafo e apresentar uma sequência de frases concisas, esclarecendo: a delimitação da temática/assunto abordado; o objetivo principal da pesquisa; a questão norteadora ou hipótese(s) do trabalho, bem como o viés teórico-metodológico que o fundamentou. Deve constar também o método e os procedimentos técnicos utilizados na obtenção e na análise dos dados; apontando principais resultados e pontos de discussão. O resumo não deve ser uma enumeração de tópicos, nem conter citações. Recomenda-se isar o verbo na voz ativa e na terceira pessoa do singular. O uso do verbo na primeira pessoa depende do diálogo com o(a) orientador(a) e o Programa de Pós-graudação no qual o Trabalho Acadêmico-Científico está sendo desenvolvido. \textbf{Seu texto deve ser digitado com espaçamento entre linhas simples}\footnote{O próprio Guia de Normalização contradiz isso. E a NBR 14724 não cita o resumo como um dos casos de espaçamento simples.} e fonte de tamanho 12. Evitar o uso de fórmulas, diagramas e siglas sem a respectiva designação. O resumo deve conter de 150 a 500 palavras e ser elaborado em conformidade com a NBR 6028:2021. Abaixo do resumo, deve-se deixar uma linha em branco e colocar a expressão ``Palavras-chave'', seguida de dois pontos e um espaço. Em seguida, deve-se informar as palavras-chave, especificando de três a cinco delas (constituída por termos extraídos de Thesaurus da área ou expressões significativas ao texto, preferencialmente que não sejam coincidentes com as palavras usadas no título do trabalho para favorecer sua busca por mecanismos de recuperação). Exceto quando ela se constituir como um substantivo próprio ou nome científico, cada palavra-chave é iniciada por letra minúscula, seguida de ponto e vírgula, e de um espaço, sendo a última finalizada por um ponto]. 

\vspace{\baselineskip} % Fazer o espacinho entre as duas partes

\noindent \textbf{Palavras-chave:} palavra-chave 1; palavra-chave 2; palavra-chave 3.

% Recomendação do Guia de Normalização de iniciar cada elemento pré-textual e textual no anverso da página
\OnesideTwoside{\clearpage}{\cleardoublepage}
  % É exatamente o mesmo código do 
% Só fico na dúvida, será que o abstract tem que ser escrito em itálico, já que é todo em outro idioma? Se sim, é só englobar todas as partes com texto (sem o \noindent) em um \textit{}
\chapter*{Abstract}


% O \noindent é necessário por causa que, como eu configurei para indentar o primeiro parágrafo no texto, como exigido, o LaTeX iria normalmente indentar o parágrafo
\noindent \lipsum[1] % Só escrever o teu abstractzinho, sem estresse

\vspace{\baselineskip} % Fazer o espacinho entre as duas partes

\noindent \textbf{Keywords:} keyword 1; keyword 2; keyword 3.

% Recomendação do Guia de Normalização de iniciar cada elemento pré-textual e textual no anverso da página
\OnesideTwoside{\clearpage}{\cleardoublepage}

  \listoffigures
  \listoftables
  \listofquadros

  % Crio um elemento utilizado para listar cada abreviatura e sigla
% Obs.: Dependendo do tamanho das siglas utilizadas, você pode querer aumentar ou diminuir o espaço entre a sigla e o seu significado, porém, mantendo o alinhamento, para isto, você modifica o comando \newcommand\siglaLargura{10ex}, substituindo o 10ex pelo tamanho que você deseja
% O tamanho 10ex suporta uma sigla de até mais ou menos 10 letras, visto que a unidade utiliza como referência o tamanho da letra "x". Você também pode utilizar a unidade "em", que utiliza como referência a letra "m", se quiser.
% Fonte: https://tex.stackexchange.com/a/508154
\newcommand\siglaLargura{10ex} % Comprimento do espaço onde fica a sigla
\newcommand\siglaGap{1ex} % Comprimento do vão entre sigla e seu significado, para dar uma folguinha no tamanho máximo
\newcommand\nomeSiglaLargura{\dimexpr\linewidth-\siglaLargura-\siglaGap\relax}
\newcommand\sigla[2]{\noindent\parbox[t]{\siglaLargura}{#1\strut}%
  \hspace{\siglaGap}%
  \parbox[t]{\nomeSiglaLargura}{#2\strut}}

% Adicione as suas siglas:
\chapter*{Lista de abreviaturas e siglas}
% Para adicionar uma abreviatura ou sigla na lista apenas precisa adicionar a sigla dentro do \sigla{<insira a sigla>}{<insira o nome completo da sigla>}

\sigla{IBGE}{Instituto Brasileiro de Geografia e Estatística}

\sigla{IFFar}{Instituto Federal de Educação, Ciência e Tecnologia Farroupilha}

\sigla{TCC}{Trabalho de Conclusão de Curso}

\sigla{T}{Teste}

\sigla{xxxxxxxxxx}{Lorem ipsum dolor sit amet, consectetur adipiscing elit. Curabitur et neque metus. Ut volutpat sem quis lacus rutrum, a viverra odio tincidunt}
  % Crio um elemento utilizado para listar cada símbolo e seu significado
% Obs.: É a mesma lógica da lista de siglas. Ajuste o tamanho do comprimento da "caixinha" do símbolo conforme a sua necessidade
% Fonte: https://tex.stackexchange.com/a/508154
\newcommand\simboloLargura{5ex} % Comprimento do espaço onde fica o símbolo
\newcommand\simboloGap{1ex} % Comprimento do vão entre símbolo e seu nome, para dar uma folguinha no tamanho máximo
\newcommand\simboloNomeLargura{\dimexpr\linewidth-\simboloLargura-\simboloGap\relax}
\newcommand\simbolo[2]{\noindent\parbox[t]{\simboloLargura}{#1\strut}%
  \hspace{\simboloGap}%
  \parbox[t]{\simboloNomeLargura}{#2\strut}}

% Adicione os seus símbolos:
\chapter*{Lista de símbolos}

% Para adicionar uma abreviatura ou sigla na lista apenas precisa adicionar o símbolo no comando \simbolo{<insira o símbolo>}{<insira o significado do símbolo>}

\simbolo{\LaTeX}{LaTeX}

\simbolo{\S}{Parágrafo}

\simbolo{$\Sigma$}{Somatória}

\simbolo{$\sum_{n = 1}^{\infty}$}{Matemágica}

\simbolo{xxxxx}{Lorem ipsum dolor sit amet, consectetur adipiscing elit. Curabitur et neque metus. Ut volutpat sem quis lacus rutrum, a viverra odio tincidunt}

% Recomendação do Guia de Normalização de iniciar cada elemento pré-textual e textual no anverso da página
\OnesideTwoside{\clearpage}{\cleardoublepage}

  \tableofcontents
\end{SecoesNaoNumeradas}
% \clearpage

% Seção ilustrativa sobre seções e subseções
% \chapter{Teste}
\begingroup
  \chapter{Exemplo de seção primária} 
  \lipsum[1]

  \section{Exemplo de seção secundária}
  \lipsum[2-4]

  \subsection{Exemplo de seção terciária}
  \lipsum[5]

  \subsubsection{Exemplo de seção quaternária}
  \lipsum[6]

  \paragraph{Exemplo de seção quinária} % Utiliza-se paragraph para o nível hierarquico inferior à subsubsection. Contudo, foi necessário customizá-lo para apresentar corretamente como um nível de seção
  \lipsum[7]

  \paragraph{Exemplo de título com indicação numérica que, ao ocupar mais de uma linha, deve ser, a partir da segunda linha, alinhado abaixo da primeira letra da primeira palavra do título}
  \lipsum[8]
\endgroup

\begingroup
  \chapter{Exemplos de elementos comuns}
\section{Alíneas e subalíneas}
  \lipsum[1-1]
  \begin{alinea}
    \item a matéria da alínea começa por letra minúscula, exceto quando se tratar de     substantivos próprios, e termina em ponto e vírgula, com exceção da última, que termina em ponto final; 
    \item o trecho final da seção correspondente, anterior às alíneas, termina em
    dois pontos;
    \item asdasdsada: 
    \begin{subalinea}
      \item a matéria da alínea começa por letra minúscula, exceto quando se tratar de     substantivos próprios, e termina em ponto e vírgula, com exceção da última, que termina em ponto final; 
      \item o trecho final da seção correspondente, anterior às alíneas, termina em
      dois pontos;
    \end{subalinea}
    \item sdsdsdaa
  \end{alinea}

  \subsection{Outros exemplos de listas numeradas}
    % \begin{leiParagrafo}
    %   \item a matéria da alínea começa por letra minúscula, exceto quando se tratar de     substantivos próprios, e termina em ponto e vírgula, com exceção da última, que termina em ponto final; 
    %   \item o trecho final da seção correspondente, anterior às alíneas, termina em
    %   dois pontos;
    %   \item asdasdsada
    %   \item sdsdsdaa
    % \end{leiParagrafo}

\section{Figuras, tabelas e quadros}
  --- Falar que figure, table e quadro são apenas invólucros, que qualquer coisa pode ser adicionada dentro, que apenas indica que o elemento possuirá seu respectivo tratamento diferenciado na hora de listar e apresentar como tal
  
\subsection{Figuras}
  --- Falar os formatos de imagem suportados
  \begin{figure}[H]
    \centering\singlespacing
    \caption{As dimensões da comunicação organizacional}
    \label{figura:marca-iffar}
    % Podem ser utilizados vários argumentos para adicionar a figura em seu devido tamanho, como: scale; width; ou height. No caso de scale, a escala, adiciona o valor em decimal representando a porcentagem da escala, já width ou height utiliza-se uma unidade de medida, como cm, px ou pt.
    \includegraphics[scale=0.35]{marca-iffar.png}\\
    \footnotesize
    Fonte: \textcite{iffar-identidade-visual-2021}.
  \end{figure}
  
\subsection{Tabelas}
\begin{table}[H]
  \centering\singlespacing

  \caption{Posições do IFFar no \textit{ranking} do painel de dados da LAI}
  \label{tabela:cgu-iffar}
  \begin{tabularx}{12cm}
    {X r r} % Definição das colunas
    \hline
    \multicolumn{1}{c}{\textbf{Categoria}}    &   
    \multicolumn{1}{c}{\textbf{Posição}}  & 
    \multicolumn{1}{c}{\textbf{Informação adicional}} \\
    \hline

    Transparência ativa &
    252\textordmasculine{} &
    24,49\% em cumprimento \\
    %\hline

    Pedidos recebidos &
    192\textordmasculine{} &
    807 pedidos \\
    %\hline

    Tempo médio de resposta para pedidos de informação       &   
    154\textordmasculine{} &
    14,91 dias \\
    \hline

  \end{tabularx}

\hspace{\fill}

\footnotesize

Fonte: Gaba guba.
\end{table}

\subsection{Quadros}
--- Falar que é uma tabela com bordinha nos lados
\begin{quadro}[H]
  \centering\singlespacing

  \caption{Posições do IFFar no \textit{ranking} do painel de dados da LAI}
  \label{quadro:cgu-iffar}
  \begin{tabularx}{12cm}
    {|X|r|r|} % Definição das colunas
    \hline
    \multicolumn{1}{|c|}{\textbf{Categoria}}    &   
    \multicolumn{1}{c}{\textbf{Posição}}  & 
    \multicolumn{1}{|c|}{\textbf{Informação adicional}} \\
    \hline

    Transparência ativa &
    252\textordmasculine{} &
    24,49\% em cumprimento \\
    %\hline

    Pedidos recebidos &
    192\textordmasculine{} &
    807 pedidos \\
    %\hline

    Tempo médio de resposta para pedidos de informação       &   
    154\textordmasculine{} &
    14,91 dias \\
    \hline

  \end{tabularx}

\hspace{\fill}

\footnotesize

Fonte: Gaba guba.
\end{quadro}
    

  Conforme o quadro no Apêndice \ref{apendice:teste} e o Anexo \ref{anexo:teste}. Citando \textcite{brasil-lei-11.892-2008}
  
  Testesddsdfdfs

\endgroup

%% Aqui são adicionados as referências, apêndices e anexos. É essencial que eles estejam dentro desse ambiente, para formatar corretamente o título da seção deles
\begin{SecoesNaoNumeradas}
%% Aqui é impressa a lista de referências bibliográficas
\begingroup
  \begin{FlushLeft} % As referências devem apresentar justificação à esquerda
    % heading=bibintoc imprime o nome da seção das referências sem numeração
    % title=Referências modifica o nome da seção, que normalmente seria "Bibliografia", por conta que é o padrão utilizado pelos documentclass book e report
    \printbibliography[heading=bibintoc, title=Referências]
  \end{FlushLeft}
\endgroup

%% Depois das referências, apresentam-se os apêndices e, depois, os anexos
\begingroup
\apendice{Um exemplo de apêndice}\label{apendice:teste}
\lipsum[6-7]
\endgroup

\begingroup 
\anexo{Um exemplo de anexo}\label{anexo:teste}
\lipsum[6-7]
\endgroup
\end{SecoesNaoNumeradas}

\end{document}